% Use the existing owner-derived source mechanism; do not edit frozen originals.
\OLSADeclareDocumentTextCorrection{OLP-0024}{FA0024-SUCCESSOR}
{تابع پسین}
{تابع جانشین}
\OLSADeclareDocumentTextCorrection{OLP-0024}{FA0024-UNDO-INTUITION}
{کاری را که $f$ انجام می‌دهد «خنثی می‌کند».}
{نتیجهٔ کار $f$ را به ورودی اولیه برمی‌گرداند.}
\OLSADeclareDocumentTextCorrection{OLP-0024}{FA0024-INVERSE-CANON-SYNONYM}
{یک \emph{وارون} برای تابع}
{یک \emph{وارون (معکوس)} برای تابع}
\OLSADeclareDocumentTextCorrection{OLP-0024}{FA0024-NATURAL-CANDIDATE}
{اکنون تعیین می‌کنیم که توابع چه هنگام وارون دارند. نامزدی مناسب
برای وارون $f\colon A \to B$ تابع $g\colon B \to A$ است که «تعریف
می‌شود» به‌وسیلهٔ}
{اکنون بررسی می‌کنیم توابع چه هنگام وارون دارند. یک پیشنهاد طبیعی
برای وارون $f\colon A \to B$ تابع $g\colon B \to A$ است که آن را با عبارت زیر
«تعریف می‌کنیم»:}
\OLSADeclareDocumentTextCorrection{OLP-0024}{FA0024-DEFINITION-CAUTION}
{اما گیومه‌های احتیاطی پیرامون «تعریف می‌شود» (و «آن») نشان می‌دهند که
این یک تعریف نیست. دست‌کم، در کلیت کامل همیشه کار
نخواهد کرد. زیرا برای آنکه این تعریف تابعی را مشخص کند، باید یگانه~$x$ای
وجود داشته باشد که $f(x) =
y$---خروجی~$g$ باید به‌طور یکتا مشخص شود.}
{اما گیومه‌های دور «تعریف می‌کنیم» (و «آن») هشدار می‌دهند که
هنوز تعریفی معتبر در همهٔ حالت‌ها نداریم. این عبارت همیشه
تابعی را مشخص نمی‌کند: باید دقیقاً یک~$x$
وجود داشته باشد که $f(x) =
y$---خروجی~$g$ باید به‌طور یکتا مشخص شود.}
\OLSADeclareDocumentTextCorrection{OLP-0024}{FA0024-NO-X-SUFFIX}
{اگر هیچ~$x$ای}
{اگر هیچ~$x$}
\OLSADeclareDocumentTextCorrection{OLP-0024}{FA0024-LEFT-UNDO}
{چنین $g$ای کاری را که $f$ انجام می‌دهد «خنثی می‌کند» و}
{تابع $g$ پس از اعمال $f$، ورودی اولیه را بازمی‌گرداند و}
\OLSADeclareDocumentTextCorrection{OLP-0024}{FA0024-RIGHT-FUNCTION-NOUN}
{چنین $h$ای}
{تابع $h$}
\OLSADeclareDocumentTextCorrection{OLP-0024}{FA0024-RIGHT-UNDO}
{$f$ کاری را که $h$ انجام می‌دهد «خنثی می‌کند».}
{$f$ پس از اعمال $h$، ورودی اولیه را بازمی‌گرداند.}
\OLSADeclareDocumentTextCorrection{OLP-0024}{FA0024-LEFT-NONEMPTY-HYPOTHESIS}
{اگر $f\colon A \to B$ !!{injective} باشد، آنگاه \emph{وارون
چپی} مانند~$g\colon B \to A$ برای~$f$ وجود دارد چنان‌که $g(f(x)) = x$ برای همهٔ $x
\in A$.}
{اگر $f\colon A \to B$ !!{injective} باشد و دامنهٔ آن ناتهی باشد، آنگاه \emph{وارون
چپی} مانند~$g\colon B \to A$ برای~$f$ وجود دارد چنان‌که $g(f(x)) = x$ برای همهٔ $x
\in A$.

\emph{یادداشت ویراستاری:} شرط ناتهی‌بودن دامنه در متن مبدأ نیامده است.
اگر دامنه تهی و هم‌دامنه ناتهی باشد، تابع تهی یک‌به‌یک است،
اما هیچ تابعی از هم‌دامنه به دامنه وجود ندارد؛ بنابراین وارون چپ ندارد.
اگر هر دو مجموعه تهی باشند، تابع تهی وارون خودش است.}
\OLSADeclareDocumentTextCorrection{OLP-0024}{FA0024-LEFT-PROOF-HYPOTHESIS}
{فرض کنید $f\colon A \to B$ !!{injective} باشد. یک $y \in B$ را در نظر بگیرید.}
{فرض کنید $f\colon A \to B$ !!{injective} باشد و دامنهٔ آن ناتهی باشد. یک $y \in B$ را در نظر بگیرید.}
\OLSADeclareDocumentTextCorrection{OLP-0024}{FA0024-UNIQUE-PREIMAGE-GRAMMAR}
{فقط یک چنین~$x \in A$ای وجود دارد.}
{تنها یک~$x \in A$ با این ویژگی وجود دارد.}
\OLSADeclareDocumentTextCorrection{OLP-0024}{FA0024-DEFINITE-X-GRAMMAR}
{$g(y)$ «آن» $x \in A$ای است که}
{$g(y)$ همان $x \in A$ است که}
\OLSADeclareDocumentTextCorrection{OLP-0024}{FA0024-FIXED-DEFAULT}
{می‌توانیم آن را به هر~$a \in A$ای نگاشت کنیم.}
{می‌توانیم آن را به یک عضو دلخواه~$a \in A$ نگاشت کنیم.}
\OLSADeclareDocumentTextCorrection{OLP-0024}{FA0024-RANGE-QUANTIFIER-GRAMMAR}
{برای هر چنین $y \in \ran{f}$ای}
{برای هر $y \in \ran{f}$}
\OLSADeclareDocumentTextCorrection{OLP-0024}{FA0024-CHOICE-NOT-IGNORED}
{در اینجا از آن چشم می‌پوشیم}
{در اینجا به‌تفصیل به آن نمی‌پردازیم}
\OLSADeclareDocumentTextCorrection{OLP-0024}{FA0024-INVERSE-ON-RANGE}
{راهی اندکی کلی‌تر برای استخراج وارون‌ها وجود دارد. در
\olref[kin]{sec} دیدیم که هر تابع $f$، !!a{surjection} $f'
\colon A \to \ran{f}$ را با قرار دادن $f'(x) = f(x)$ برای همهٔ $x \in A$ القا می‌کند.}
{روش اندکی کلی‌تری برای به دست آوردن وارون وجود دارد. در
\olref[kin]{sec} دیدیم که از هر تابع $f$ می‌توان یک !!a{surjection} $f'
\colon A \to \ran{f}$ را با قرار دادن $f'(x) = f(x)$ برای همهٔ $x \in A$ به دست آورد.}
\OLSADeclareDocumentTextCorrection{OLP-0025}{FA0025-INVERSE-GRAMMAR}
{\oliflabeldef{sfr:fun:inv:sec}{در \olref[inv]{sec} دیدیم که
وارونِ~$f^{-1}$ِ !!a{bijection}~$f$ خود تابع است. عمل دیگری
که روی توابع انجام می‌دهیم، ترکیب است. اکنون می‌توانیم}{می‌توانیم}}
{\oliflabeldef{sfr:fun:inv:sec}{در \olref[inv]{sec} دیدیم که
$f^{-1}$، یعنی وارونِ !!a{bijection}~$f$، خود یک تابع است. عمل دیگری
که روی توابع انجام می‌دهیم، ترکیب است. اکنون می‌توانیم}{می‌توانیم}}
\OLSADeclareDocumentTextCorrection{OLP-0025}{FA0025-ASSIGNMENT-CONSTRUCTION}
{$(\comp{f}{g}) \colon A \to C$ هر !!{element} از~$A$
را با !!a{element}ی از~$C$ جفت می‌کند. اینکه کدام !!{element} از~$C$
با !!a{element}ی از $A$ جفت می‌شود، چنین مشخص می‌گردد: با داشتن ورودی $x \in
A$، نخست تابع $f$ را بر~$x$ اعمال می‌کنیم که مقداری مانند $f(x)
= y \in B$ خروجی می‌دهد؛ سپس تابع $g$ را بر~$y$ اعمال می‌کنیم که مقداری مانند
$g(f(x)) = g(y) = z \in C$ خروجی می‌دهد.}
{$(\comp{f}{g}) \colon A \to C$ به هر !!{element} از~$A$
!!a{element}ی از~$C$ نسبت می‌دهد. برای مشخص کردن اینکه کدام !!{element} از~$C$
به !!a{element}ی از $A$ نسبت داده می‌شود، این روش را به کار می‌بریم: با داشتن ورودی $x \in
A$، نخست تابع $f$ را بر~$x$ اعمال می‌کنیم که مقدار $f(x)
= y \in B$ را به دست می‌دهد؛ سپس تابع $g$ را بر~$y$ اعمال می‌کنیم که مقدار
$g(f(x)) = g(y) = z \in C$ را به دست می‌دهد.}
\OLSADeclareDocumentTextCorrection{OLP-0025}{FA0025-CAPTION}
{\caption{ترکیبِ $g \circ f$ از دو تابع $f$ و~$g$.}}
{\caption{$g \circ f$، ترکیب دو تابع $f$ و~$g$.}}
\OLSADeclareDocumentTextCorrection{OLP-0025}{FA0025-SUCCESSOR}
{پسین آن را بگیرید}
{جانشین آن را بگیرید}
\OLSADeclareDocumentTextCorrection{OLP-0025}{FA0025-GRAPH}
{نشان دهید گراف}
{نشان دهید نمودار}
\OLSADeclareDocumentTextCorrection{OLP-0026}{FA0026-GRAPH-DEFINITION}
{\begin{defn}[گراف یک تابع جزئی]
فرض کنید $f\colon A \pto B$ تابعی جزئی باشد. \emph{گرافِ}~$f$}
{\begin{defn}[نمودار یک تابع جزئی]
فرض کنید $f\colon A \pto B$ تابعی جزئی باشد. \emph{نمودارِ}~$f$}
\OLSADeclareDocumentTextCorrection{OLP-0026}{FA0026-GRAPH-PROPOSITION}
{پس $R$ گراف تابع جزئی}
{پس $R$ نمودار تابع جزئی}
\OLSADeclareSetup{OLP-0024}{\olsa_source_correction_install_document_text:}
\OLSADeclareSetup{OLP-0025}{\olsa_source_correction_install_document_text:}
\OLSADeclareSetup{OLP-0026}{\olsa_source_correction_install_document_text:}
